\section{Cronograma}

As atividades serão desenvolvidas de forma incremental ao longo do restante do semestre 2026.2, abrangendo desde o estudo das técnicas e a preparação dos dados até a integração, avaliação e apresentação do sistema. O cronograma a seguir apresenta os períodos previstos para cada etapa. Todas as datas referem-se ao ano de 2026.

\begin{table}[H]
    \centering
    \caption{Cronograma previsto para o desenvolvimento do projeto.}
    \label{tab:cronograma}
    \begin{tabularx}{\textwidth}{p{3cm} X}
        \toprule
        \textbf{Período} & \textbf{Atividade} \\
        \midrule

        21/09 a 27/09 &
        Estudo da literatura, análise das alternativas de detecção e rastreamento, definição da metodologia e preparação inicial do ambiente de desenvolvimento. \\

        28/09 a 11/10 &
        Obtenção e organização dos dados, análise das informações disponíveis, leitura dos vídeos e implementação do fluxo inicial de processamento quadro a quadro. \\

        12/10 a 25/10 &
        Desenvolvimento e integração das etapas de detecção e rastreamento dos jogadores e investigação de métodos para identificação das equipes a partir de características visuais. \\

        26/10 a 08/11 &
        Implementação e validação da transformação das coordenadas dos jogadores do plano da imagem para o plano do campo utilizando homografia. \\

        09/11 a 22/11 &
        Integração das etapas do sistema, construção das trajetórias dos jogadores, desenvolvimento da visualização dos resultados, implementação da interface experimental e realização dos testes de avaliação. \\

        23/11 a 02/12 &
        Comparação dos resultados obtidos com as informações de referência, análise das limitações do método, consolidação dos resultados, elaboração do relatório final e preparação da apresentação. \\

        03/12 a 10/12 &
        Apresentação do projeto, realização de eventuais ajustes finais e entrega. \\

        \bottomrule
    \end{tabularx}
\end{table}